\scanpage{26}
1) groupe Alg. $\boxed{G}$ sur $\mathbb Q_p$. (Il dépend en fait du choix de $\Omega$. cf plus bas 2\textsuperscript{o})

2) Hom \(\boxed{\rho:\pi\longrightarrow G(\mathbb Q_p)}\)

[ L’op. de $\pi=\operatorname{Aut}_L(\Omega)$ sur $G=G^{(\Omega)}$ est donnée par $\gamma\to\operatorname{int}(\rho(\gamma))$ ]

\deleted{3)} On fait opérer $\pi$ sur $G_\Omega$ \uline{diagonalement}, soit $G_\Omega^{!}$

3) Hom \(\boxed{j_{\deleted{\unsure{s}}}^{!}:G_{m_\Omega}\longrightarrow G_\Omega^{!}}\) \qquad inv. par $\pi$
\reversemarginpar\margin{$G_\Omega$\quad \unsure{1|!}}\normalmarginpar

On appelle $\deleted{\unsure{T_p}}\inserted{S_\Omega}(1)$ \deleted{\unsure{$T$}} le \deleted{\unsure{rep}} $\Omega$-module inversible $T_p(\Omega)^{\unsure{\times|\ast}}\otimes_{\mathbb Z_p}\Omega$ avec op. diagonale de $\pi$, \deleted{\unsure{associé}} et $\unsure{\mathcal T|\tau}_\Omega$ le $G_{m_\Omega}$ torseur associé. Posons

\deleted{$\mathcal T\mathbin{\wedge}_{\unsure{G_m}}^{\unsure{\Phi|\mathfrak{P}}}G_{m_\Omega}\quad G_{m_{\unsure{p}}}(\Omega)\quad\unsure{Q_\Omega}\quad\mathcal T\mathbin{\wedge}_{G_{m_\Omega}}\unsure{\Phi|\mathfrak{P}}_\Omega$}
\inserted{$Q_\Omega=\mathcal T\mathbin{\wedge}\unsure{\Phi|P|\mathfrak{P}}$} torseur \inserted{à dr} sous $G_\Omega$ à droite, à gauche est un \uline{$G_\Omega^{!}$ tordu} par $\mathcal T$, soit \tikz[baseline=(x.base)]{\node[draw,ellipse,inner sep=1pt](x){$G_\Omega^{\prime\prime}$};}.

[ correspondant au foncteur fibre \inserted{(gradué)} sur $\operatorname{Rep}(G)$
\[
G_{m_\Omega}\xrightarrow{j_\Omega^{\unsure{\prime\prime|\ast}}}G_\Omega^{\prime\prime}
\qquad M\to\sum\operatorname{gr}^{i}_{j}(M_\Omega)\otimes S_\Omega(-i)=\mathrm{Hdg}_\Omega(M)
\]

\unsure{$\subset$} \uline{Axiome} Le \deleted{foncteur}\inserted{torseur} $Q_\Omega$ \deleted{\unsure{sous}} se descend à $L$ en un $G_L$-torseur $Q_L$

\begin{center}\tikz[baseline=(x.base)]{\node[draw,ellipse,inner sep=3pt](x){$Q_L\quad G_L$};}\end{center}

(unique à isom unique près)

en un $G_L^{\prime\prime}$, et
\[
G_{m_L}\xrightarrow{j_L}G_L^{\prime\prime}
\]
Correspond à fct fibre gradué $\mathrm{Hdg}_L$

\margin{\uline{NB} Un $j$ satisfaisant ces conditions est unique. Signifie que les $M_\Omega^{\deleted{\unsure{0}}}$ \uline{sont de Hodge} Du moins si}

\noindent\rule{\linewidth}{0.4pt}

\textcircled{\unsure{1|a|2\textsuperscript{o}}} Il y a deux foncteurs fibres de $\operatorname{Rep}(G)$ dans $\operatorname{Mod}(\Omega,\pi)$

a) $M\longrightarrow M_\Omega$ opérations \unsure{provenant|par transport} des op. ordinaires de $\pi$ sur $\Omega$

b) $M\unsure{\longrightarrow|\mathbin{-}|\mapsto}M_\Omega^{!}$ opérations diagonales

- foncteurs $T_{\unsure{p\Omega|G\Omega|p,\Omega}}$ et $T_{p\Omega}^{!}$

\deleted{Après oubli des op.} On \unsure{a} alors un \deleted{\unsure{$(\Omega,\pi)$-}}torseur \inserted{à droite} sur $(\Omega,\pi)$ du groupe $G_\Omega$, soit $\Phi_\Omega$, donné par le 1-cocycle $\rho:\pi\unsure{\longrightarrow|\mathbin{-}}G(\mathbb Q_p)\subset H^0(\Omega,\pi\deleted{\unsure{G}},G_\Omega)=G(\unsure{L|k|\Omega})$. Le \deleted{groupe}\inserted{$(\Omega,\pi)$-} $\underline{\operatorname{Aut}}_{G_\Omega,\pi}(\Phi)$ est $G_\Omega^{!}$.
% archive sheet number: 25
